\documentclass{subfiles}
\begin{document}
    Differently from the two algorithms previously described, 
        Lenstra's factorization algorithm assumes we are working with an elliptic curve \(C*\)
        over the ring \(\Z / n\Z\).
        The key idea is that, while computing \(kP\), for some \(k \in \Z\) and \(P \in \C*\),
        we will eventually have to compute the inverse of a element of the ring that admtis no inverse.
        Precisely we proceed as follow:
        \begin{enumerate}
            \item Let \(\C*\) be the elliptic curve \(Y^{2} = X^{3} + AX + B\)
                over the ring \(\Z/n\Z\).

            \item Fix a point \(P = (a, b)\) at random and \(B \equiv b^{2} - a^{3} - aA \mod{n}\).

            \item For \(j = 2, 3, 4, \ldots\)  up to some bound do:
                \begin{enumerate}
                    \item Compute \(Q \equiv jP \mod{N}\) and set \(P = Q\).
                    \item If the previous step fails, we have found a non trivial divisor \(d\).
                    \item If \(d \le n\), return \(d\).
                    \item If \(d = n\) go back a choose a new point, a new curve or both.
                \end{enumerate}
            \item Go back to step 2.
        \end{enumerate}

    Lenstra's factorization algorithm requires \(\bigO{e^{\sqrt{(\log n)(\log\log n)}}}\)
    time.
\end{document}
